\documentclass{situationist-preprint}

\title{Relativism Refuted by Meta-Disagreement}
\msid{MTE-06-a}
\msfamily{MTE-06}
\msdiscipline{Metaethics}
\mstype{Principal manuscript}
\mskeywords{moral relativism; meta-disagreement; self-refutation; faultless disagreement; peritrope; objectivity}

\begin{document}
\maketitle

\begin{abstract}
Moral relativism holds that moral claims are true only relative to frameworks, and that when frameworks disagree neither party is mistaken. We argue that the view is self-refuting, not at the level of first-order disagreement, where it is consistent, but at the level of meta-disagreement: disagreement about whether a first-order disagreement is one in which someone is mistaken. Every actual moral tradition contains such meta-claims, since its members hold that those who reject its central verdicts are in error. The relativist must say whether her own verdict on a meta-disagreement, that neither party is mistaken, is true simpliciter or only relative to her framework. If it is true simpliciter, then, given the disquotational schema for truth, it entails both the claim in dispute and its negation. If it is true only relative to her framework, then by her own account the objectivist's contrary verdict is true relative to the objectivist's framework, and relativism entails that it is false of that framework. Either way relativism refutes itself. The argument does not tell against the nature-relativity of moral content defended in MTE-02.
\end{abstract}

\section{The problem}

Moral relativism is the view that moral claims are not true or false simpliciter but only relative to a moral framework, the standards of a person or a community, and that there is no framework-independent fact that makes one framework correct and another mistaken. In its classical form the view is Protagoras's, as Plato reports it \citep[152a]{Theaetetus}. Its modern defenders have given it more careful formulations: \citet{Harman1975} relativizes moral judgments to implicit agreements between speaker and agent, \citet{Wong1984,Wong2006} to the adequate moralities of communities, and recent semantic relativists relativize the truth of evaluative assertions to contexts of assessment, so that two thinkers can disagree without either being at fault \citep{Kolbel2004,MacFarlane2014}.

The oldest objection is that relativism refutes itself. Plato's version, the \emph{peritrope}, observes that if every belief is true for the believer, then the belief that Protagoras's doctrine is false is true for those who hold it \citep[170e--171c]{Theaetetus}. \citet{Burnyeat1976} reconstructed the argument and showed that it works against a relativism of truth stated without restriction. Moral relativism, however, is restricted to the moral domain, and its defenders have replied that the doctrine itself, being a claim about morality rather than a moral claim, falls outside its own scope. The thesis is then consistent: it is true simpliciter that moral claims are only relatively true.

We argue that the reply fails, and that the self-refutation reappears one level up. The relativist's characteristic verdict is not about first-order moral claims but about disagreements over them: that when two frameworks disagree, neither party is mistaken. Moral traditions issue the contrary verdict. They hold that those who reject their central claims are in error. The disagreement between these verdicts is a \emph{meta-disagreement}, and relativism cannot coherently take a side in it.

\section{Framework}

\begin{definition}\label{d:framework}
A \emph{moral framework} is a body of moral claims accepted by a person or community, together with the claims about error that its members draw from them. We write $\mathrm{T}_F(p)$ when $p$ is true relative to framework $F$.
\end{definition}

\begin{definition}\label{d:rel}
\emph{Relativism} is the conjunction of two theses. \emph{Relativity}: for every moral claim $p$ and framework $F$, $p$ is true relative to $F$ just in case $F$ supports it. \emph{Faultlessness}: if $F$ supports $p$ and $G$ supports $\neg p$, then neither the members of $F$ nor the members of $G$ are mistaken about $p$.
\end{definition}

\begin{definition}\label{d:meta}
Two frameworks \emph{meta-disagree} about a first-order disagreement over $p$ if one of them supports the claim that the other's members are mistaken about $p$, and the other supports the claim that neither is.
\end{definition}

\begin{assumption}[Objectivist frameworks]\label{a:objfw}
Some moral framework $F$ supports a first-order claim $p$ and also supports the claim that anyone whose framework supports $\neg p$ is mistaken about $p$.
\end{assumption}

\begin{assumption}[Disquotation]\label{a:tschema}
Truth simpliciter satisfies the disquotational schema: the judgment that $q$ is true simpliciter if and only if $q$. A person who judges that $q$ is mistaken about it if and only if her judgment is not true simpliciter.
\end{assumption}

Assumption~\ref{a:objfw} is a claim about actual moral traditions, and it is not in doubt. Those who hold that slavery is wrong hold that slaveholders who judged it permissible were mistaken, not that they were right by their own lights; the claim belongs to the framework, as \citet{Dworkin1996} argued, and is not a philosophical gloss added to it. Empirical work on folk metaethics confirms that ordinary moral judges treat many moral disagreements as disagreements in which one party is wrong \citep{GoodwinDarley2008}. Assumption~\ref{a:tschema} states the minimal content of truth simpliciter, which the relativist cannot deny without abandoning the contrast between relative truth and truth that her view relies on.

\section{The argument}

Let $F$ be as in Assumption~\ref{a:objfw}, and let $G$ be a framework, such as the relativist's own, that supports $\neg p$ or, more weakly, supports faultlessness for the disagreement over $p$ between $F$ and a framework that supports $\neg p$. Then $F$ and the relativist meta-disagree. The relativist's verdict on the meta-disagreement is the faultlessness thesis applied to it: neither party to the first-order disagreement is mistaken. Call this verdict $N$.

\begin{lemma}\label{l:abs}
If $N$ is true simpliciter and the first-order disagreement is between a framework supporting $p$ and one supporting $\neg p$, then $p$ and $\neg p$ are both true.
\end{lemma}

\begin{proof}
If the members of $F$, who judge that $p$, are not mistaken, then by Assumption~\ref{a:tschema} their judgment is true simpliciter, and so $p$. If the members of the other framework, who judge that $\neg p$, are not mistaken, their judgment is true simpliciter, and so $\neg p$.
\end{proof}

\begin{lemma}\label{l:rel}
If $N$ is true only relative to frameworks, then relativism, applied to the framework $F$, is false relative to $F$.
\end{lemma}

\begin{proof}
By Relativity, the claim that the other framework's members are mistaken about $p$ is true relative to $F$, since $F$ supports it (Assumption~\ref{a:objfw}). Faultlessness, applied to the disagreement over $p$, entails that they are not mistaken. Since $N$ is only relatively true, its negation is true relative to $F$, and so is the negation of faultlessness for this case. Relativism therefore entails that, relative to $F$, relativism is false.
\end{proof}

\begin{theorem}\label{t:main}
Relativism is self-refuting at the level of meta-disagreement: it is either inconsistent or false by its own account of truth for every objectivist framework.
\end{theorem}

\begin{proof}
The relativist's verdict $N$ is either true simpliciter or true only relative to frameworks. In the first case relativism, which asserts $N$, entails a contradiction by Lemma~\ref{l:abs}. In the second, by Lemma~\ref{l:rel}, relativism entails that it is false relative to every framework that satisfies Assumption~\ref{a:objfw}; since such frameworks exist, relativism entails that, by the only standard of truth it recognizes for those frameworks, it is false of them. A thesis that claims to describe the status of every framework's moral claims, and that by its own lights is false of some frameworks, is refuted by its own content.
\end{proof}

The argument does not require that the relativist assert anything about first-order moral claims simpliciter. It requires only that she give a verdict on meta-disagreements, which her view exists to give: faultlessness is what distinguishes relativism from the view that moral disagreement is difficult. The reply that relativism is a claim about morality rather than a moral claim, and so outside its own scope, does not help. If the claim that neither party is mistaken is outside the scope of relativity, it is true simpliciter, and the first horn applies; if it is within it, the second does.

\section{Objections}

\begin{objection}[Rigidified relativism]
A sophisticated relativist holds that \enquote{wrong}, in each framework, refers rigidly to what that framework's actual standards condemn. Then the members of $F$ speak truly when they say that those who approve of slavery are mistaken, and the relativist agrees. There is no meta-disagreement.
\end{objection}

\begin{reply}
If the relativist agrees, without qualification, that those who approve of slavery are mistaken, she has abandoned faultlessness and is no longer a relativist. If she agrees only that \enquote{they are mistaken} is true as uttered within $F$, while holding that the slaveholders' contrary judgment is equally true as uttered within theirs, she has restated faultlessness, and the members of $F$, who intend their claim of error to hold whatever the slaveholders' standards, reject it. The rigidified semantics changes where the meta-disagreement is located, not whether it exists.
\end{reply}

\begin{objection}[Assessment sensitivity]
On the semantics of \citet{MacFarlane2014}, the truth of an assertion is relative to a context of assessment, and the norms of retraction and disagreement are defined accordingly. Faultless disagreement is coherent there, and no contradiction arises.
\end{objection}

\begin{reply}
The semantics is consistent, and on it $N$ is itself assessment-sensitive; that is the second horn. An assessor in the context of $F$ correctly assesses the claim that neither party is mistaken as false, and is required by the semantics to reject it. The relativist who uses the semantics to defend relativism to the members of $F$ is defending a claim that, by the semantics she relies on, they are correct to reject.
\end{reply}

\begin{objection}[Objectivist frameworks rest on an error]
The members of $F$ believe that their claims of error hold simpliciter, and they are wrong; the relativist is entitled to correct them.
\end{objection}

\begin{reply}
To correct them, the relativist must hold that their meta-claim is false simpliciter, and so that some claim belonging to a moral framework, the claim that the slaveholders are mistaken, has a truth value independent of frameworks. That is the first horn, and by Lemma~\ref{l:abs} it brings with it the first-order claim and its negation.
\end{reply}

\section{What the result does not establish}

The theorem refutes relativism about the truth of moral claims. It does not establish that every moral claim has a determinate truth value, and it is compatible with the view that some moral questions have no answer. It does not refute the view that moral \emph{content} is relative to the nature of the agents concerned. MTE-02 argues that morality is objective but not absolute: what is required of an agent depends on her nature, but every rational agent, judging under idealized conditions, reaches the same verdict on every moral claim, including claims about agents of other natures. That view involves no faultless disagreement, since it holds each moral claim true or false simpliciter, and so it generates no meta-disagreement of the kind exploited here. Nor does the theorem show that moral disagreement is evidence against realism, a question addressed in MTE-12.

\section{Conclusion}

Relativism is consistent about first-order moral claims. Its distinctive verdict, that disagreeing frameworks are faultless, is a verdict on meta-disagreements, and every moral tradition that holds its opponents to be in error issues the contrary verdict. If the relativist's verdict is true simpliciter, it entails a contradiction; if it is true only relative to frameworks, relativism is false relative to the objectivist's framework by its own standard. Relativism is therefore self-refuting at the level of meta-disagreement.

\begin{thebibliography}{99}\small
\bibitem[Burnyeat(1976)]{Burnyeat1976} Burnyeat, M.~F. (1976). Protagoras and self-refutation in Plato's \emph{Theaetetus}. \emph{The Philosophical Review}, 85(2), 172--195.
\bibitem[Dworkin(1996)]{Dworkin1996} Dworkin, R. (1996). Objectivity and truth: You'd better believe it. \emph{Philosophy \& Public Affairs}, 25(2), 87--139.
\bibitem[Goodwin and Darley(2008)]{GoodwinDarley2008} Goodwin, G.~P. and Darley, J.~M. (2008). The psychology of meta-ethics: Exploring objectivism. \emph{Cognition}, 106(3), 1339--1366.
\bibitem[Harman(1975)]{Harman1975} Harman, G. (1975). Moral relativism defended. \emph{The Philosophical Review}, 84(1), 3--22.
\bibitem[K\"olbel(2004)]{Kolbel2004} K\"olbel, M. (2004). Faultless disagreement. \emph{Proceedings of the Aristotelian Society}, 104, 53--73.
\bibitem[MacFarlane(2014)]{MacFarlane2014} MacFarlane, J. (2014). \emph{Assessment Sensitivity: Relative Truth and Its Applications}. Oxford University Press.
\bibitem[Plato()]{Theaetetus} Plato. \emph{Theaetetus}.
\bibitem[Wong(1984)]{Wong1984} Wong, D.~B. (1984). \emph{Moral Relativity}. University of California Press.
\bibitem[Wong(2006)]{Wong2006} Wong, D.~B. (2006). \emph{Natural Moralities: A Defense of Pluralistic Relativism}. Oxford University Press.
\end{thebibliography}

\end{document}
